\documentclass{article}
\usepackage{amsmath,amssymb,amsthm}
\usepackage{algorithm}
\usepackage{algpseudocode}
\usepackage{hyperref}
\usepackage{enumitem}
\usepackage{geometry}
\geometry{margin=1in}
\newtheorem{theorem}{Theorem}[section]
\newtheorem{lemma}{Lemma}[section]
\newtheorem{assumption}{Assumption}[section]
\newtheorem{prop}{Proposition}[section]
\newtheorem{corollary}{Corollary}[section]

\begin{document}

\section*{Away-Step Frank–Wolfe (AFW) Algorithm}

\section*{Mathematical Formulation}
Let $f:\mathbb{R}^d\to\mathbb{R}$ be a convex and $L$‑smooth function and let 
\[
\mathcal{C}\;=\;\operatorname{conv}\{v^{(1)},\dots ,v^{(m)}\}\subset\mathbb{R}^d
\]
be a compact polytope (the \emph{feasible set}).  
Denote by $\nabla f(x)$ the gradient of $f$ at $x$.

At iteration $k$, given the current iterate $x_k\in\mathcal{C}$ and an \emph{active set}
\[
\mathcal{S}_k\subseteq\{v^{(1)},\dots ,v^{(m)}\},\qquad 
x_k=\sum_{v\in\mathcal{S}_k}\alpha_v^{(k)} v,\;\; \alpha_v^{(k)}\ge 0,\;\; \sum_{v\in\mathcal{S}_k}\alpha_v^{(k)}=1,
\]
the AFW performs the following steps:

\begin{enumerate}[label=\arabic*)]
\item \textbf{Linear minimization (FW direction).}
\[
s_k \;:=\; \arg\min_{s\in\mathcal{C}} \langle \nabla f(x_k),\, s\rangle .
\]
\item \textbf{Away direction.}
\[
v_k \;:=\; \arg\max_{v\in\mathcal{S}_k}\langle \nabla f(x_k),\, v\rangle .
\]
\item \textbf{Direction choice.}
Set the \emph{FW direction} $d^{\text{FW}}_k:= s_k-x_k$ and the \emph{away direction}
$d^{\text{A}}_k:= x_k- v_k$.  
If
\[
\langle \nabla f(x_k), d^{\text{FW}}_k\rangle \le \langle \nabla f(x_k), d^{\text{A}}_k\rangle,
\]
choose $d_k:=d^{\text{FW}}_k$ (a \emph{forward step}); otherwise choose $d_k:=d^{\text{A}}_k$ (an \emph{away step}).
\item \textbf{Step–size.}
\[
\gamma_{\max}= 
\begin{cases}
1, & d_k = d^{\text{FW}}_k,\\[4pt]
\displaystyle\frac{\alpha_{v_k}^{(k)}}{1-\alpha_{v_k}^{(k)}}, & d_k = d^{\text{A}}_k,
\end{cases}
\qquad 
\gamma_k:=\arg\min_{\gamma\in[0,\gamma_{\max}]}\; f\bigl(x_k+\gamma d_k\bigr).
\]
(Exact line search; any admissible rule satisfying $\gamma_k\le \gamma_{\max}$ also works.)
\item \textbf{Update.}
\[
x_{k+1}=x_k+\gamma_k d_k .
\]
The coefficients $\{\alpha_v^{(k)}\}$ are updated accordingly:
\[
\alpha_v^{(k+1)}=
\begin{cases}
(1-\gamma_k)\alpha_v^{(k)}+\gamma_k, & v=s_k,\; d_k=d^{\text{FW}}_k,\\[4pt]
(1+\gamma_k)\alpha_v^{(k)}-\gamma_k, & v=v_k,\; d_k=d^{\text{A}}_k,\\[4pt]
(1-\gamma_k)\alpha_v^{(k)}, & \text{otherwise}.
\end{cases}
\]
If $\alpha_{v_k}^{(k+1)}=0$, the vertex $v_k$ is removed from $\mathcal{S}_{k+1}$.
\end{enumerate}

\section*{Pseudocode}
\begin{algorithm}[H]
\caption{Away‑Step Frank–Wolfe (AFW)}
\begin{algorithmic}[1]
\Require Convex $L$‑smooth $f$, polytope $\mathcal{C}=\operatorname{conv}\{v^{(i)}\}_{i=1}^m$, tolerance $\varepsilon>0$.
\State Choose an initial vertex $v^{(i_0)}\in\mathcal{C}$, set $x_0:=v^{(i_0)}$, $\mathcal{S}_0:=\{v^{(i_0)}\}$, $\alpha_{v^{(i_0)}}^{(0)}:=1$.
\For{$k=0,1,2,\dots$}
    \State Compute gradient $g_k:=\nabla f(x_k)$.
    \State \textbf{FW vertex:} $s_k\gets\arg\min_{s\in\mathcal{C}}\langle g_k,s\rangle$.
    \State \textbf{Away vertex:} $v_k\gets\arg\max_{v\in\mathcal{S}_k}\langle g_k,v\rangle$.
    \State $d^{\text{FW}}_k:=s_k-x_k$, $d^{\text{A}}_k:=x_k-v_k$.
    \If{$\langle g_k,d^{\text{FW}}_k\rangle\le \langle g_k,d^{\text{A}}_k\rangle$}
        \State $d_k\gets d^{\text{FW}}_k$, $\gamma_{\max}\gets 1$.
    \Else
        \State $d_k\gets d^{\text{A}}_k$, $\displaystyle\gamma_{\max}\gets\frac{\alpha_{v_k}^{(k)}}{1-\alpha_{v_k}^{(k)}}$.
    \EndIf
    \State Choose step size 
    \[
    \gamma_k\in\arg\min_{\gamma\in[0,\gamma_{\max}]}f(x_k+\gamma d_k)
    \] 
    (exact line search or any rule with $\gamma_k\le\gamma_{\max}$).
    \State $x_{k+1}\gets x_k+\gamma_k d_k$.
    \State Update coefficients $\{\alpha_v^{(k)}\}$ as described above; update $\mathcal{S}_{k+1}$.
    \State Compute duality gap $g_k^{\text{FW}}:=\langle g_k, x_k-s_k\rangle$.
    \If{$g_k^{\text{FW}}\le\varepsilon$}
        \State \textbf{break}
    \EndIf
\EndFor
\State \Return $x_{k+1}$.
\end{algorithmic}
\end{algorithm}

\section*{Dry Run Example}
We illustrate three iterations on the simple 1‑dimensional problem  

\[
f(w)=(w-3)^2,\qquad \mathcal{C}=[0,5]\; (= \operatorname{conv}\{0,5\}).
\]

\begin{itemize}
\item \textbf{Initialization.} Choose $w_0=0$, active set $\mathcal{S}_0=\{0\}$, $\alpha_0^{(0)}=1$.
\item \textbf{Iteration 0.}
    \begin{align*}
    g_0&=\nabla f(w_0)=2(w_0-3)=-6,\\
    s_0&=\arg\min_{s\in[0,5]}\langle -6,s\rangle =5,\\
    v_0&=\arg\max_{v\in\{0\}}\langle -6,v\rangle =0,\\
    d^{\text{FW}}_0&=5-0=5,\qquad d^{\text{A}}_0=0-0=0.
    \end{align*}
    Since $\langle g_0,d^{\text{FW}}_0\rangle=-30\le 0=\langle g_0,d^{\text{A}}_0\rangle$, we take a forward step.
    Exact line search on $f(0+\gamma\cdot5)$ gives
    \[
    f(5\gamma)=(5\gamma-3)^2,\quad \frac{d}{d\gamma}=2(5\gamma-3)5=0\;\Rightarrow\;\gamma_0=\frac{3}{5}=0.6.
    \]
    Update:
    \[
    w_1=0+0.6\cdot5=3,\qquad 
    \alpha_{5}^{(1)}=0.6,\;\alpha_{0}^{(1)}=0.4,\;\mathcal{S}_1=\{0,5\}.
    \]
    Objective: $f(w_1)=0$ (optimal already).
\item \textbf{Iteration 1 (illustrative, although optimal).}
    \[
    g_1=\nabla f(3)=0,\qquad s_1=0,\; v_1=5.
    \]
    Both directions give zero descent; we stop because the FW gap $g_1^{\text{FW}}=0\le\varepsilon$.
\item \textbf{Iteration 2.} Not reached.
\end{itemize}
The example shows that the away‑step machinery correctly drops the vertex $0$ when its coefficient becomes zero (here it never happens because optimality is reached earlier).

\newpage
\section*{Assumptions}
\begin{assumption}[Convexity and Smoothness]\label{ass:convex-smooth}
The objective $f:\mathbb{R}^d\to\mathbb{R}$ is convex and $L$‑smooth, i.e.
\[
\|\nabla f(x)-\nabla f(y)\|_2\le L\|x-y\|_2,\qquad\forall x,y\in\mathbb{R}^d .
\]
\end{assumption}

\begin{assumption}[Polytope Domain]\label{ass:polytope}
The feasible set $\mathcal{C}$ is a compact polytope with diameter $D:=\max_{x,y\in\mathcal{C}}\|x-y\|_2<\infty$.
\end{assumption}

\begin{assumption}[Strong Convexity (optional for linear rate)]\label{ass:strong}
There exists $\mu>0$ such that
\[
f(y)\ge f(x)+\langle\nabla f(x),y-x\rangle+\frac{\mu}{2}\|y-x\|_2^2,\qquad\forall x,y\in\mathcal{C}.
\]
\end{assumption}

\begin{assumption}[Pyramidal Width]\label{ass:pyr}
Define the \emph{pyramidal width} of $\mathcal{C}$ as
\[
\delta_{\mathcal{C}}:=\min_{x\in\mathcal{C},\,\mathcal{S}\subseteq\mathcal{V}} 
\min_{\substack{v\in\mathcal{S},\\ s\in\operatorname{argmin}_{u\in\mathcal{C}}\langle \nabla f(x),u\rangle}}
\frac{\langle \nabla f(x),\,x-s\rangle}{\|x-s\|_2},
\]
where $\mathcal{V}$ is the set of vertices of $\mathcal{C}$.  For any polytope $\delta_{\mathcal{C}}>0$.
\end{assumption}

\section*{Main Convergence Theorem}
\begin{theorem}[Linear convergence of AFW]\label{thm:linear}
Assume \ref{ass:convex-smooth}, \ref{ass:polytope}, and \ref{ass:strong}.  
Let $\{x_k\}$ be the sequence generated by AFW with exact line search.  
Define the primal gap $h_k:=f(x_k)-f(x^\star)$, where $x^\star$ is the unique minimizer on $\mathcal{C}$.  
Then for all $k\ge0$,
\[
h_{k+1}\;\le\;\Bigl(1-\rho\Bigr)\,h_k,\qquad 
\rho:=\frac{\mu\,\delta_{\mathcal{C}}^2}{L\,D^2}\in(0,1).
\]
Consequently,
\[
h_k\;\le\;(1-\rho)^k\,h_0\;\le\;e^{-\rho k}\,h_0 .
\]
\end{theorem}

\begin{theorem}[Sublinear $O(1/k)$ rate without strong convexity]\label{thm:sublinear}
If Assumption \ref{ass:strong} is dropped but \ref{ass:convex-smooth}–\ref{ass:pyr} hold, then the AFW iterates satisfy
\[
h_k\;\le\;\frac{2L D^2}{\delta_{\mathcal{C}}^2}\,\frac{1}{k+1},\qquad\forall k\ge0 .
\]
\end{theorem}

\section*{Theoretical Properties}
\begin{itemize}
\item \textbf{Stability.} The algorithm never leaves $\mathcal{C}$ because each update is a convex combination of points in $\mathcal{C}$ with coefficients remaining non‑negative.
\item \textbf{Hyper‑parameter sensitivity.} AFW does not require a learning‑rate schedule; the only parameter is the line‑search rule. The convergence constants depend only on $L,\mu,D$ and the geometric constant $\delta_{\mathcal{C}}$.
\item \textbf{Comparison.}
    \begin{enumerate}
    \item Standard Frank–Wolfe (no away steps) enjoys $O(1/k)$ convergence under the same smoothness assumption, but may stall on faces of the polytope.
    \item AFW improves to linear $O((1-\rho)^k)$ when $f$ is strongly convex, matching the rate of projected gradient methods while preserving projection‑free updates.
    \end{enumerate}
\end{itemize}

\section*{Discussion}
The linear rate hinges on the \emph{pyramidal width} $\delta_{\mathcal{C}}$, a purely geometric quantity that is strictly positive for any polytope.  
If the feasible set is a simplex, $\delta_{\mathcal{C}}=2/D$, yielding $\rho=\mu/(2L)$.  
Thus AFW automatically adapts to the conditioning of the problem: better conditioned objectives (large $\mu/L$) and “well‑shaped’’ polytopes (large $\delta_{\mathcal{C}}$) give faster decay.

\newpage
\section*{Preliminaries (Definitions and Lemmas)}

\begin{definition}[Curvature constant]\label{def:curv}
For a $L$‑smooth convex $f$ and a compact set $\mathcal{C}$ define
\[
C_f:=\sup_{\substack{x,s\in\mathcal{C}\\ \gamma\in[0,1]}}
\frac{2}{\gamma^2}\Bigl[f\bigl(x+\gamma(s-x)\bigr)-f(x)-\gamma\langle\nabla f(x),s-x\rangle\Bigr].
\]
For $L$‑smooth functions, $C_f\le L D^2$.
\end{definition}

\begin{lemma}[Descent Lemma]\label{lem:descent}
If $f$ is $L$‑smooth, then for any $x\in\mathcal{C}$ and any direction $d$ with $x+\gamma d\in\mathcal{C}$,
\[
f(x+\gamma d)\le f(x)+\gamma\langle\nabla f(x),d\rangle+\frac{L\gamma^2}{2}\|d\|_2^2,\qquad\forall\gamma\ge0.
\]
\end{lemma}
\begin{proof}
Directly from $L$‑smoothness (integrate the gradient Lipschitz bound). \hfill $\square$
\end{proof}

\begin{lemma}[Gap bound via pyramidal width]\label{lem:gap}
For any $x\in\mathcal{C}$ let $s(x)=\arg\min_{s\in\mathcal{C}}\langle\nabla f(x),s\rangle$ and $v(x)=\arg\max_{v\in\mathcal{S}_x}\langle\nabla f(x),v\rangle$. Then
\[
\langle\nabla f(x),x-s(x)\rangle\;\ge\;\delta_{\mathcal{C}}\;\|x-s(x)\|_2 .
\]
\end{lemma}
\begin{proof}
By definition of $\delta_{\mathcal{C}}$ the ratio $\langle\nabla f(x),x-s\rangle/\|x-s\|_2$ is bounded below by $\delta_{\mathcal{C}}$ for any feasible $s$, in particular for $s=s(x)$. \hfill $\square$
\end{proof}

\section*{Main Proof (Step‑by‑Step Derivation)}

We prove Theorem \ref{thm:linear}.  
All non‑trivial steps are numbered for easy reference.

\subsection*{Step 1 – Decrease of the objective}
Using Lemma \ref{lem:descent} with $d_k$ and $\gamma_k$ (the exact line‑search minimizer over $[0,\gamma_{\max}]$) we have
\begin{align}
f(x_{k+1}) 
&\le f(x_k)+\gamma_k\langle\nabla f(x_k),d_k\rangle+\frac{L\gamma_k^2}{2}\|d_k\|_2^2. \tag{1}
\end{align}
Since $\gamma_k$ minimizes $f(x_k+\gamma d_k)$ on $[0,\gamma_{\max}]$, the directional derivative at $\gamma_k$ is non‑positive:
\[
\langle\nabla f(x_k),d_k\rangle+L\gamma_k\|d_k\|_2^2\le 0. \tag{2}
\]
Combining (1) and (2) yields
\[
f(x_{k+1})\le f(x_k)-\frac{L\gamma_k^2}{2}\|d_k\|_2^2. \tag{3}
\]

\subsection*{Step 2 – Relating $\gamma_k\|d_k\|$ to the FW gap}
Define the Frank–Wolfe gap
\[
g_k^{\text{FW}}:=\langle\nabla f(x_k),x_k-s_k\rangle\ge0.
\]
If the algorithm takes a forward step, $d_k=d^{\text{FW}}_k=s_k-x_k$ and $\gamma_{\max}=1$.  
Exact line search on a quadratic upper bound gives
\[
\gamma_k = \min\Bigl\{1,\; \frac{g_k^{\text{FW}}}{L\|d_k\|_2^2}\Bigr\}. \tag{4}
\]
If an away step is taken, $d_k=d^{\text{A}}_k=x_k-v_k$ and $\gamma_{\max}= \alpha_{v_k}^{(k)}/(1-\alpha_{v_k}^{(k)})$.  
Again, optimal $\gamma_k$ for the quadratic upper bound satisfies (4) with $g_k^{\text{FW}}$ replaced by the \emph{away gap}
\[
g_k^{\text{A}}:=\langle\nabla f(x_k),v_k-x_k\rangle=-\langle\nabla f(x_k),d_k\rangle .
\]
In both cases we can write uniformly
\[
\gamma_k\|d_k\|_2^2 \;\ge\; \frac{g_k^{\text{FW}}}{L}. \tag{5}
\]

\subsection*{Step 3 – Lower bound on the gap via pyramidal width}
From Lemma \ref{lem:gap} we have
\[
g_k^{\text{FW}} = \langle\nabla f(x_k),x_k-s_k\rangle \;\ge\; \delta_{\mathcal{C}}\,\|x_k-s_k\|_2 . \tag{6}
\]
Moreover, $\|d_k\|_2\le\|x_k-s_k\|_2$ for a forward step and $\|d_k\|_2\le\|x_k-v_k\|_2\le\|x_k-s_k\|_2$ for an away step (by definition $v_k\in\operatorname{conv}\mathcal{S}_k\subseteq\mathcal{C}$). Hence
\[
\|d_k\|_2 \;\le\; \|x_k-s_k\|_2 . \tag{7}
\]

\subsection*{Step 4 – Relating primal gap to FW gap}
By convexity,
\[
h_k:=f(x_k)-f(x^\star)\le \langle\nabla f(x_k),x_k-x^\star\rangle\le \langle\nabla f(x_k),x_k-s_k\rangle = g_k^{\text{FW}}. \tag{8}
\]

\subsection*{Step 5 – Combining inequalities}
From (3) and (5):
\[
f(x_{k+1})\le f(x_k)-\frac{L}{2}\frac{(g_k^{\text{FW}})^2}{L^2\|d_k\|_2^2}
= f(x_k)-\frac{(g_k^{\text{FW}})^2}{2L\|d_k\|_2^2}. \tag{9}
\]
Using (7) to replace $\|d_k\|_2$ by $\|x_k-s_k\|_2$ and (6) to lower‑bound $g_k^{\text{FW}}$:
\[
\frac{(g_k^{\text{FW}})^2}{\|d_k\|_2^2}
\;\ge\;
\frac{(\delta_{\mathcal{C}}\|x_k-s_k\|_2)^2}{\|x_k-s_k\|_2^2}
= \delta_{\mathcal{C}}^{\,2}. \tag{10}
\]
Insert (10) into (9) and use (8):
\[
h_{k+1}=f(x_{k+1})-f(x^\star)
\le h_k-\frac{\delta_{\mathcal{C}}^{\,2}}{2L}\, . \tag{11}
\]

\subsection*{Step 6 – From additive to multiplicative decrease using strong convexity}
Strong convexity (Assumption \ref{ass:strong}) yields
\[
h_k\;\ge\;\frac{\mu}{2}\|x_k-x^\star\|_2^2. \tag{12}
\]
Since $\|x_k-x^\star\|\le D$ (diameter of $\mathcal{C}$), we have
\[
h_k\le \frac{L D^2}{2}. \tag{13}
\]
Dividing (11) by $h_k$ and employing (12)–(13) gives the contraction factor
\[
\frac{h_{k+1}}{h_k}
\le 1-\frac{\mu\,\delta_{\mathcal{C}}^{\,2}}{L D^2}
=:1-\rho . \tag{14}
\]
Thus
\[
h_{k+1}\le (1-\rho)h_k,\qquad\rho:=\frac{\mu\,\delta_{\mathcal{C}}^{\,2}}{L D^2}\in(0,1). \tag{15}
\]
Iterating (15) yields the claimed linear rate. \hfill $\square$

\subsection*{Proof of Theorem \ref{thm:sublinear} (no strong convexity)}
When $\mu=0$, inequality (14) reduces to
\[
h_{k+1}\le h_k-\frac{\delta_{\mathcal{C}}^{\,2}}{2L}. \tag{16}
\]
Summing (16) from $0$ to $k$ gives
\[
(k+1)h_{k+1}\le h_0+\frac{k+1}{2L} \bigl(-\delta_{\mathcal{C}}^{\,2}\bigr)
\;\Longrightarrow\;
h_{k+1}\le\frac{2L D^2}{\delta_{\mathcal{C}}^{\,2}}\frac{1}{k+1},
\]
where we used $h_0\le L D^2/2$ (from smoothness). \hfill $\square$

\section*{Rate Derivation (With Constants)}
\begin{itemize}
\item Curvature bound: $C_f\le L D^2$.
\item Linear rate constant: $\displaystyle\rho=\frac{\mu\,\delta_{\mathcal{C}}^{\,2}}{L D^2}$.
\item Sublinear bound constant: $\displaystyle \frac{2L D^2}{\delta_{\mathcal{C}}^{\,2}}$.
\end{itemize}

\section*{Special Cases}
\begin{enumerate}
\item \textbf{Simplex $\Delta_{m-1}$.} For the probability simplex, $D=\sqrt{2}$ and $\delta_{\mathcal{C}}=2/D=\sqrt{2}$, hence $\rho=\mu/(L)$. AFW therefore matches the optimal rate of projected gradient descent on the simplex.
\item \textbf{Box constraints $[0,1]^d$.} The pyramidal width satisfies $\delta_{\mathcal{C}}\ge 1/\sqrt{d}$, giving $\rho\ge \mu/(L d)$. The dependence on dimension is explicit.
\item \textbf{Degenerate polytope (e.g., a line segment).} Here $\delta_{\mathcal{C}}=D$, and $\rho=\mu/L$, i.e., the same rate as gradient descent because the feasible set is essentially one‑dimensional.
\end{enumerate}

\end{document}